In answer to our questions, our \robot\ can tell us which direction it is facing.
Since it cannot speak, its answers are {\em boolean}, \index{Boolean} that is 
{\mcl True} ({\em true} \index{True}) or {\mcl False} 
({\em false} \index{False}), written in \python\ with a capital letter. We can imagine that it gives its answers with an indicator light
which it switches on to answer  {\mcl True} and leaves off to answer {\mcl False}.

To ask questions about the orientation of the \robot, we use the following test commands:

\begin{table}[!tb]
\begin{center}
\begin{tabular}{@{}|l|p{11.7cm}|@{}} \hline
\multicolumn{2}{|l|}{\em Commands testing the orientation of the \robot:} \\ \hline \hline
{\idt tn()} & {\mcl True} if  the \robot\ is \underline{t}urned towards the north (\underline{n}) and {\mcl False} otherwise. \\ \hline
{\idt tne()} & {\mcl True} if  the \robot\ is \underline{t}urned towards the north-east (\underline{ne}) and {\mcl False} otherwise. \\ \hline
{\idt te()} & {\mcl True} if  the \robot\ is \underline{t}urned towards the east (\underline{e}) and {\mcl False} otherwise. \\ \hline
{\idt tse()} & {\mcl True} if  the \robot\ is \underline{t}urned towards the south-east (\underline{se}) and {\mcl False} otherwise. \\ \hline
{\idt ts()} & {\mcl True} if  the \robot\ is \underline{t}urned towards the south (\underline{s}) and {\mcl False} otherwise. \\ \hline
{\idt tsw()} & {\mcl True} if  the \robot\ is \underline{t}urned towards the south-west (\underline{sw}) and {\mcl False} otherwise. \\ \hline
{\idt tw()} & {\mcl True} if  the \robot\ is \underline{t}urned towards the west (\underline{w}) and {\mcl False} otherwise. \\ \hline
{\idt tnw()} & {\mcl True} if  the \robot\ is \underline{t}urned towards the north-west (\underline{nw}) and {\mcl False} otherwise. \\ \hline
\end{tabular}
\end{center}
%\label{tab:tnetc}
\caption{Tests of the orientation of the \protect\robot}
\end{table}

In \python, the ``if'' statement makes it possible to execute a sequence of instructions only
if a certain {\em boolean expression} \index{Boolean expression}is true ({\mcl True}).  To begin with, we will
use simple boolean expressions, namely the commands testing the orientation of the \robot\
({\idt tn()}, {\idt tne()}, {\idt te()}, {\idt tse()}, {\idt ts()}, {\idt tsw()}, {\idt tw()}, {\idt tnw()}). In its
simplest form, the ``if'' statement is written as follows:

\index{Test}
\begin{center}
\begin{tabular}{@{}|p{12.6cm}|@{}} \hline
{\em ``if'' statement with a single alternative:}\\ \hline \hline
\begin{minipage}[t]{12.6cm}
{\medskip
\mbox{\blc\ \ \ {\mcl if}\hspace{0.3em}\blanc\hspace{0.3em}{\small\rm boolean expression}%
\footnote{\ With result {\mcl True} or {\mcl False}.}{\idt :}}\par 
\mbox{\blc\ \ \ \ \ \ \ \small\rm  \ldots\ Instructions (indented) executed only if}\par 
\mbox{\blc\ \ \ \ \ \ \ \small\rm  the boolean expression is {\mcl True} \ldots}\par 
\medskip}
\end{minipage}\\ \hline 
\end{tabular}
\end{center}

It reads ``if \ldots\ then''.
As for the body of a procedure or of a {\mcl for} loop, the line of the {\mcl if} ends with
a colon `{\idt :}' and the instructions that depend on the test are indented (shifted four
spaces to the right with respect to the {\mcl if}): they form a {\em block} which ends at the first line
that is less indented. If we forget the colon, \python\ rejects the program with the message
`{\tt SyntaxError: expected ':'}' and if we forget to indent,
with the message
`{\tt IndentationError: expected an indented block after 'if' statement on line 5}'
(if the {\mcl if} is on line 5 of the program).

When there is only one instruction to execute, it can be written on the same line, after the
colon, as in `{\mcl if}\ {\idt tw():\ qtr()}'. But beware: in
`{\mcl if}\ {\idt tw():\ qtr();\ fd()}' both instructions {\idt qtr()} and {\idt fd()} depend on the test.
\underline{A serious error} is to get the indentation wrong. For example in:
{
\parindent \parindentprog
\medskip

\vspace{\neglskip}\mbox{\blc {\mcl if}{\idt \ tw():}}\par
\vspace{\neglskip}\mbox{\blc \blk\blk{\idt qtr()}}\par
\vspace{\neglskip}\mbox{\blc {\idt fd()}}\par

\medskip }
\noindent only the instruction {\idt qtr()} depends on the test: the instruction {\idt fd()},
which is not indented, is always ex\-e\-cut\-ed, whether the \robot\ is facing west or not.
\python\ cannot report this error since the program is correct (but does not do
what we wanted). Note also that the ``if'' statement is {\em one} instruction:
if other instructions follow it, they are written on the following lines, with the same
indentation as the {\mcl if}.


In the example below, the \robot\ moves forward twice and turns a quarter turn to the right, then
one more quarter turn 
if it is facing west:

\vspace{\baselineskip}

 {\raggedright 
 \includevspacePostScript[0.9cm]{FriseEnPeigne}{0.95}{0.95}
 \input python/CombFrieze.tex
}
 
 In its first form above, the ``if'' statement does nothing when the boolean expression
 is {\mcl False}. If we want to choose the sequence of instructions to execute in both
 alternatives, we use the following form:
 
\index{Test}
 \begin{center}
\begin{tabular}{@{}|p{12.6cm}|@{}} \hline
{\em ``if'' statement with two alternatives:}\\ \hline \hline
\begin{minipage}[t]{12.6cm}
{\medskip
\mbox{\blc\ \ \ {\mcl if}\hspace{0.3em}\blanc\hspace{0.3em}{\small\rm boolean expression}%
\footnote{ \ With result {\mcl True} or {\mcl False}.}{\idt :}}\par 
\mbox{\blc\ \ \ \ \ \ \ \small\rm  \ldots\ Instructions (indented) executed if the boolean}\par 
\mbox{\blc\ \ \ \ \ \ \ \small\rm   expression is {\mcl True} \ldots}\par 
\mbox{\blc\ \ \ {\mcl else}{\idt :}}\par 
\mbox{\blc\ \ \ \ \ \ \ \small\rm  \ldots\ Instructions  (indented) executed if the boolean  }\par 
\mbox{\blc\ \ \ \ \ \ \ \small\rm  expression is {\mcl False} \ldots}\par 
\medskip}
\end{minipage}\\ \hline 
\end{tabular}
\end{center}

It reads ``if \ldots\ then \ldots\ otherwise \ldots'' ({\mcl else} means `otherwise').

The {\mcl else} is written exactly below the {\mcl if} (with the same indentation) and it is followed
by a colon. As before, an alternative consisting of a single instruction
can be written on the same line as the {\mcl if} or the {\mcl else}, after the colon.
Forgetting the colon after the {\mcl else} or misaligning the
{\mcl else} below the {\mcl if} is an error, but not a very
serious one since it is reported by \python\ before execution (`{\tt SyntaxError: expected ':'}'
or `{\tt IndentationError: unindent does not match any outer indentation level}').
\python\ also reports an error (`{\tt SyntaxError: invalid syntax}') when we forget
to indent one of the instructions placed between the {\mcl if} and the {\mcl else}.

In the drawing of the sawtooth frieze below, the \robot\ turns three quarters of a turn to the right 
if it is facing north-east and to the left in the other cases (the only other possible case being, in fact,
that it is facing south):

\vspace{0.4cm}
 {\raggedright 
 \includevspacePostScript[0.5cm]{FriseEnDentsDeScie}{0.95}{0.95}
 \input python/SawtoothFrieze.tex
 }
 \vfill
 When there are more than two cases to test, we can use several alternatives and do nothing
 when none of the boolean expressions is {\mcl True}, as in:
 
\index{Test}
 \begin{center}
\begin{tabular}{@{}|p{12.6cm}|@{}} \hline
{\em ``if'' statement with $n$ alternatives:}\\ \hline \hline
\begin{minipage}[t]{12.6cm}
{\medskip
\vspace{\neglskip}\mbox{\blc\ \ \ {\mcl if}\hspace{0.3em}\blanc\hspace{0.3em}{\small\rm boolean expression 1}{\idt :}}\par 
\vspace{\neglskip}\mbox{\blc\ \ \ \ \ \ \ \small\rm  \ldots\ Instructions (indented) executed if boolean }\par 
\vspace{\neglskip}\mbox{\blc\ \ \ \ \ \ \ \small\rm  expression 1 is {\mcl True} \ldots}\par 
\vspace{\neglskip}\mbox{\blc\ \ \ {\mcl elif}\hspace{0.3em}\blanc\hspace{0.3em}{\small\rm boolean expression 2}{\idt :}}\par 
\vspace{\neglskip}\mbox{\blc\ \ \ \ \ \ \ \small\rm  \ldots\ Instructions (indented) executed if boolean }\par 
\vspace{\neglskip}\mbox{\blc\ \ \ \ \ \ \ \small\rm  expression 2 is {\mcl True} \ldots}\par 
\mbox{\blc\ \ \ \dots}\par
\vspace{\neglskip}\mbox{\blc\ \ \ {\mcl elif}\hspace{0.3em}\blanc\hspace{0.3em}{\small\rm boolean expression $n$}{\idt :}}\par 
\vspace{\neglskip}\mbox{\blc\ \ \ \ \ \ \ \small\rm  \ldots\ Instructions (indented) executed if boolean }\par 
\vspace{\neglskip}\mbox{\blc\ \ \ \ \ \ \ \small\rm  expression $n$ is {\mcl True} \ldots}\par 
\medskip}
\end{minipage}\\ \hline 
\end{tabular}
\end{center}

\vfill

 \noindent We can also provide a last alternative which is executed when all the
 boolean expressions tested are {\mcl False}:
 
 \index{Test}
\begin{center}
\begin{tabular}{@{}|p{12.6cm}|@{}} \hline
{\em ``if'' statement with $(n + 1)$ alternatives:}\\ \hline \hline
\begin{minipage}[t]{12.6cm}
{\medskip
\vspace{\neglskip}\mbox{\blc\ \ \ {\mcl if}\hspace{0.3em}\blanc\hspace{0.3em}{\small\rm boolean expression 1}{\idt :}}\par 
\vspace{\neglskip}\mbox{\blc\ \ \ \ \ \ \ \small\rm  \ldots\ Instructions (indented) executed if boolean }\par 
\vspace{\neglskip}\mbox{\blc\ \ \ \ \ \ \ \small\rm  expression 1 is {\mcl True} \ldots}\par 
\vspace{\neglskip}\mbox{\blc\ \ \ {\mcl elif}\hspace{0.3em}\blanc\hspace{0.3em}{\small\rm boolean expression 2}{\idt :}}\par 
\vspace{\neglskip}\mbox{\blc\ \ \ \ \ \ \ \small\rm  \ldots\ Instructions (indented) executed if boolean }\par 
\vspace{\neglskip}\mbox{\blc\ \ \ \ \ \ \ \small\rm  expression 2 is {\mcl True} \ldots}\par 
\mbox{\blc\ \ \ \dots}\par
\vspace{\neglskip}\mbox{\blc\ \ \ {\mcl elif}\hspace{0.3em}\blanc\hspace{0.3em}{\small\rm boolean expression $n$}{\idt :}}\par 
\vspace{\neglskip}\mbox{\blc\ \ \ \ \ \ \ \small\rm  \ldots\ Instructions (indented) executed if boolean }\par 
\vspace{\neglskip}\mbox{\blc\ \ \ \ \ \ \ \small\rm  expression $n$ is {\mcl True} \ldots}\par 
\vspace{\neglskip}\mbox{\blc\ \ \ {\mcl else}{\idt :}}\par 
\vspace{\neglskip}\mbox{\blc\ \ \ \ \ \ \ \small\rm  \ldots\ Instructions  (indented) executed if boolean expressions}\par 
\vspace{\neglskip}\mbox{\blc\ \ \ \ \ \ \ \small\rm  1, \ldots, $n$ are all {\mcl False} \ldots}\par 
\medskip}
\end{minipage}\\ \hline 
\end{tabular}
\end{center}

 The word {\mcl elif} is a contraction of {\mcl else}~{\mcl if}. The boolean
 expressions are evaluated in order and only the instructions of the first alternative
 whose expression is {\mcl True} are executed.
 This form of multiple test is illustrated below in the program that builds the following frieze:

  {\raggedright 
 \includevspacePostScript[1.0cm]{FriseDeVaguelettes}{0.95}{0.95}
 \input python/FriezeOfRipples.tex
 }

To finish, we are going to introduce new commands of the \robot\ which
make the \robot\ turn on the spot to face a given direction. For example the command
{\idt fn()}, which makes the \robot\ \underline{f}ace \underline{n}orth, can
be defined with the following \python\ procedure:

{
\parindent \parindentprog
\medskip

\vspace{\neglskip}\mbox{\blc {\cmt \#\ Turn\ the\ robot\ on\ the\ spot\ to\ face\ north}}\par
\vspace{\neglskip}\mbox{\blc {\mcl def}{\idt \ fn():}}\par
\vspace{\neglskip}\mbox{\blc \blk\blk{\mcl if}{\idt \ tne():}}\par
\vspace{\neglskip}\mbox{\blc \blk\blk\blk\blk{\idt etl()}}\par
\vspace{\neglskip}\mbox{\blc \blk\blk{\mcl elif}{\idt \ te():}}\par
\vspace{\neglskip}\mbox{\blc \blk\blk\blk\blk{\idt qtl()}}\par
\vspace{\neglskip}\mbox{\blc \blk\blk{\mcl elif}{\idt \ tse():}}\par
\vspace{\neglskip}\mbox{\blc \blk\blk\blk\blk{\idt tetl()}}\par
\vspace{\neglskip}\mbox{\blc \blk\blk{\mcl elif}{\idt \ ts():}}\par
\vspace{\neglskip}\mbox{\blc \blk\blk\blk\blk{\idt ht()}}\par
\vspace{\neglskip}\mbox{\blc \blk\blk{\mcl elif}{\idt \ tsw():}}\par
\vspace{\neglskip}\mbox{\blc \blk\blk\blk\blk{\idt tetr()}}\par
\vspace{\neglskip}\mbox{\blc \blk\blk{\mcl elif}{\idt \ tw():}}\par
\vspace{\neglskip}\mbox{\blc \blk\blk\blk\blk{\idt qtr()}}\par
\vspace{\neglskip}\mbox{\blc \blk\blk{\mcl elif}{\idt \ tnw():}}\par
\vspace{\neglskip}\mbox{\blc \blk\blk\blk\blk{\idt etr()}}\par

\medskip }

\noindent The program {\idt FacingNorth.py} of the folder {\idt book/\allowbreak 15\_if} tries this
procedure in the eight directions. We could define in the same way the commands that turn the \robot\
on the spot to face one of the eight main directions of the compass rose, which are given in table
\ref{tab:pnetc} below.

\begin{table}[htbp]
\begin{center}
\begin{tabular}{@{}|l|p{11.7cm}|@{}} \hline
\multicolumn{2}{|l|}{\em Commands turning the \robot\ on the spot:} \\ \hline \hline
{\idt fn()} & makes the \robot\ turn on the spot to \underline{f}ace north (\underline{n}). \\ \hline
{\idt fne()} &  makes the \robot\ turn on the spot to \underline{f}ace north-east (\underline{ne}). \\ \hline
{\idt fe()} &  makes the \robot\ turn on the spot to \underline{f}ace east (\underline{e}). \\ \hline
{\idt fse()} &  makes the \robot\ turn on the spot to \underline{f}ace south-east (\underline{se}). \\ \hline
{\idt fs()} &  makes the \robot\ turn on the spot to \underline{f}ace south (\underline{s}). \\ \hline
{\idt fsw()} &  makes the \robot\ turn on the spot to \underline{f}ace south-west (\underline{sw}). \\ \hline
{\idt fw()} &  makes the \robot\ turn on the spot to \underline{f}ace west (\underline{w}). \\ \hline
{\idt fnw()} &  makes the \robot\ turn on the spot to \underline{f}ace north-west (\underline{nw}). \\ \hline
\end{tabular}
\end{center}
\caption{Turning the \protect\robot\ on the spot to face a given direction\label{tab:pnetc}}
\end{table}

\needspace{6\baselineskip}\noindent \rule{137mm}{0.1mm}\par\nopagebreak %---------------------------------------------
\begin{exercice}{Write \python\ programs containing ``if''
statements to produce the friezes of figure \ref{fig:frisif} below.
\noindent $\Box$}\end{exercice}
\begin{figure}[htbp]
 \includevspacePostScript[3.8cm]{FriseIf}{0.95}{0.95}
 \caption{Friezes of circles and (half-) octagons}\label{fig:frisif}
\end{figure}
\begin{corrige}{
We give the program drawing the frieze of octagons of figure \ref{fig:frisif}
 (page \pageref{fig:frisif}), the others being in the folder {\idt book/\allowbreak 15\_if}:
\input python/FriezeOfOctagons.tex
\noindent $\Box$}\end{corrige}
\vfill\pagebreak[3]

%%%%%%%%%%%%%%%%%%%%%%%%%%%%%%%%%%%%%%%%%%%%%%%%%%%%%%%%%%%%%%%%%%%%%%%%%%%
\noindent \framebox[\textwidth]{
\begin{minipage}{\myparboxwidth}
\setlength{\parindent}{2em}
\small \sf
\vspace{\abovedisplayskip}
\section*{Addenda}
\index{Test}
The test statement {\idt if} \ldots {\idt then} \ldots {\idt else} \ldots\ was introduced towards the end of the
1950s in the programming language {\sc algol 60}. This language, designed by a group
of international experts, had a very great influence on the design of later programming languages,
notably {\sc algol w}, \pascal, {\sc ada},  {\sc modula} and many others.  
Another form of test statement, by cases, was introduced in the programming language 
{\sc algol w} (also designed by N. Wirth) and taken up in \pascal\ under the name ``case''.
In \ocaml, the test statement by cases is called {\em pattern matching} (``match'').
\python\ also has, since version 3.10, a {\idt match} statement, but we will not use it:
in \python, a test by cases is usually written as a sequence of {\mcl if} \ldots\ {\mcl elif} \ldots\
{\mcl else}, comparing the value of an expression with each of the possible cases.

For example {\idt valOrientation()} (``orientation value'') is
equal to the number of the orientation of the \robot\ (0 for north, 1 for north-east, 2 for east, 
etc.). The library {\idt LibRobot} declares the constants {\idt north}, {\idt northEast}, {\idt east},
{\idt southEast}, {\idt south}, {\idt southWest}, {\idt west} and {\idt northWest},
equal to $0$, $1$, \ldots, $7$.
We can write:
 

\index{Test}
{
\parindent \parindentprog
\medskip

\vspace{\neglskip}\mbox{\blc {\mcl if}{\idt \ valOrientation()\ ==\ north:\ \ \ \ \ \ \ \ \ }{\cmt \#\ north}}\par
\vspace{\neglskip}\mbox{\blc \blk\blk{\idt ...}}\par
\vspace{\neglskip}\mbox{\blc {\mcl elif}{\idt \ valOrientation()\ ==\ northEast:\ \ \ }{\cmt \#\ north{-}east}}\par
\vspace{\neglskip}\mbox{\blc \blk\blk{\idt ...}}\par
\vspace{\neglskip}\mbox{\blc {\mcl elif}{\idt \ valOrientation()\ ==\ east:\ \ \ \ \ \ \ \ }{\cmt \#\ east}}\par
\vspace{\neglskip}\mbox{\blc \blk\blk{\idt ...}}\par
\vspace{\neglskip}\mbox{\blc {\mcl elif}{\idt \ valOrientation()\ ==\ southEast:\ \ \ }{\cmt \#\ south{-}east}}\par
\vspace{\neglskip}\mbox{\blc \blk\blk{\idt ...}}\par
\vspace{\neglskip}\mbox{\blc {\mcl elif}{\idt \ valOrientation()\ ==\ south:\ \ \ \ \ \ \ }{\cmt \#\ south}}\par
\vspace{\neglskip}\mbox{\blc \blk\blk{\idt ...}}\par
\vspace{\neglskip}\mbox{\blc {\mcl elif}{\idt \ valOrientation()\ ==\ southWest:\ \ \ }{\cmt \#\ south{-}west}}\par
\vspace{\neglskip}\mbox{\blc \blk\blk{\idt ...}}\par
\vspace{\neglskip}\mbox{\blc {\mcl elif}{\idt \ valOrientation()\ ==\ west:\ \ \ \ \ \ \ \ }{\cmt \#\ west}}\par
\vspace{\neglskip}\mbox{\blc \blk\blk{\idt ...}}\par
\vspace{\neglskip}\mbox{\blc {\mcl elif}{\idt \ valOrientation()\ ==\ northWest:\ \ \ }{\cmt \#\ north{-}west}}\par
\vspace{\neglskip}\mbox{\blc \blk\blk{\idt ...}}\par
\vspace{\neglskip}\mbox{\blc {\mcl else}{\idt :\ \ \ \ \ \ \ \ \ \ \ \ \ \ \ \ \ \ \ \ \ \ \ \ \ \ \ \ \ \ \ \ \ }{\cmt \#\ impossible}}\par
\vspace{\neglskip}\mbox{\blc \blk\blk{\mcl pass}}\par

\medskip }

\vspace{\abovedisplayskip}
\end{minipage}
}

%%%%%%%%%%%%%%%%%%%%%%%%%%%%%%%%%%%%%%%%%%%%%%%%%%%%%%%%%%%%%%%%%%%%%%%%%%%%
%%%%%%%%%%%%%%%%%%%%%%%%%%%%%%%%%%%%%%%%%%%%%%%%%%%%%%%%%%%%%%%%%%%%%%%%%%%
\noindent \framebox[\textwidth]{
\begin{minipage}{\myparboxwidth}
\setlength{\parindent}{2em}
\small \sf
\vspace{\abovedisplayskip}
\noindent Note that the comparison is written with two equals signs `{\idt ==}', the single sign
`{\idt =}' being used to declare constants. With
`{\mcl if}\ {\idt valOrientation()\ =\ north:}', \python\ would answer
`{\tt SyntaxError: cannot assign to function call here. Maybe you meant '==' instead of '='?}'
(a function call cannot be given a value).
Two numbers can also be compared with {\idt !=} (different from, $\neq$), {\idt \char60{}}, {\idt \char60{}=} ($\leq$),
{\idt \char62{}} and {\idt \char62{}=} ($\geq$).
The last case {\mcl else} groups together all the other values, like the word
{\mcl otherwise} of some extensions of \pascal. 
The {\mcl pass} statement does nothing; this last case can also be omitted.
Several similar cases can be grouped together 
as in the example below, where we test whether the remainder of the division of the orientation number by 2 is zero:


% example of grouping cases
{
\parindent \parindentprog
\medskip

\vspace{\neglskip}\mbox{\blc {\mcl if}{\idt \ valOrientation()\ \%\ 2\ ==\ 0:\ \ }{\cmt \#\ N,\ E,\ S,\ W}}\par
\vspace{\neglskip}\mbox{\blc \blk\blk{\idt ...}}\par
\vspace{\neglskip}\mbox{\blc {\mcl else}{\idt :\ \ \ \ \ \ \ \ \ \ \ \ \ \ \ \ \ \ \ \ \ \ \ \ \ \ }{\cmt \#\ other\ cases}}\par
\vspace{\neglskip}\mbox{\blc \blk\blk{\idt ...}}\par

\medskip }

Note that with {\mcl if} tests, the cases can be names of constants such as {\idt north}
as well as numbers such as {\idt 0}, which makes the program more readable.
This was not the case with pattern matching in \ocaml, where `{\idt \char124{}\ north\ {-}\char62{}}' does not mean
`if the orientation is north' but declares a new identifier {\idt north} which matches
all values (the same is true of the {\idt match} of \python).
 
\vspace{\abovedisplayskip}
\end{minipage}
}
%%%%%%%%%%%%%%%%%%%%%%%%%%%%%%%%%%%%%%%%%%%%%%%%%%%%%%%%%%%%%%%%%%%%%%%%%%%% 
\newpage                 % If the chapter ends with a blank page
\thispagestyle{empty}
